\chapter{Expression des besoins}
\label{ch:besoins}

\projet{} étant un projet réalisé \emph{pendant la formation}, sans commanditaire externe,
c'est moi qui formule l'expression des besoins définissant les objectifs et les
limites du projet, comme le prévoit le référentiel de certification.

\section{Contexte et problème adressé}

\projet{} est un \textbf{réseau social de microblogging} inspiré de X (anciennement Twitter) :
un utilisateur publie de courts messages, éventuellement accompagnés d'une image, auxquels les
autres peuvent répondre dans un fil imbriqué et réagir par un « j'aime ».

Le besoin pédagogique sous-jacent est de se confronter à \emph{l'ensemble des briques d'une
application web sécurisée multicouche} sur un périmètre fonctionnel maîtrisable :
authentification, opérations de création / lecture / mise à jour / suppression (CRUD), relations
récursives (réponses imbriquées), téléversement de fichiers, recherche et modération. Ce
périmètre, volontairement resserré, permet de traiter \textbf{la sécurité à chaque couche}
plutôt que d'empiler des fonctionnalités superficielles.

\section{Objectifs du projet}

Les objectifs sont formulés de manière \emph{vérifiable} :

\begin{itemize}
  \item permettre à un visiteur de \textbf{s'inscrire}, de \textbf{se connecter} et de se
        déconnecter ;
  \item permettre à un utilisateur authentifié de \textbf{publier} un message (avec image
        facultative), d'y \textbf{répondre} dans un fil imbriqué, de l'\textbf{aimer}, de
        \textbf{rechercher} des utilisateurs et des messages, et d'\textbf{éditer son profil} ;
  \item permettre à un \textbf{administrateur} de \textbf{modérer} (supprimer tout contenu,
        accéder à une page d'administration) ;
  \item garantir, de manière transverse, la \textbf{sécurité applicative} (protection contre
        les failles XSS, CSRF, injection SQL, IDOR) conformément aux recommandations de l'OWASP
        et de l'ANSSI, ainsi que la conformité au RGPD et la prise en compte du RGAA.
\end{itemize}

\section{Périmètre et limites}
\label{sec:limites}

Le \textbf{hors-périmètre} est assumé explicitement : il traduit un arbitrage de maturité, et
non un oubli. Ne font pas partie du projet :

\begin{itemize}
  \item la messagerie privée entre utilisateurs ;
  \item les notifications en temps réel (WebSocket) ;
  \item la fédération inter-instances et les API publiques ouvertes ;
  \item toute fonctionnalité de paiement ou de monétisation.
\end{itemize}

\section{Acteurs et cas d'usage}
\label{sec:acteurs}

Trois acteurs interagissent avec le système, selon une relation de \emph{généralisation} :
l'administrateur est un utilisateur particulier, lui-même un visiteur authentifié.

\begin{table}[ht]
\centering
\renewcommand{\arraystretch}{1.3}
\small
\begin{tabularx}{\textwidth}{@{}L{3cm} X@{}}
\toprule
\textbf{Acteur} & \textbf{Cas d'usage principaux} \\
\midrule
Visiteur & Consulter la page de connexion, s'inscrire, se connecter \\
Utilisateur & Consulter le fil, publier, répondre, aimer, rechercher, éditer son profil, supprimer son propre contenu \\
Administrateur & Tous les cas de l'utilisateur \emph{plus} la modération (supprimer tout contenu, page d'administration) \\
\bottomrule
\end{tabularx}
\caption{Acteurs et cas d'usage.}
\label{tab:acteurs}
\end{table}

Quelques \emph{user stories} représentatives, qui serviront de base au diagramme de cas
d'utilisation (section~\ref{sec:uml}) et au plan de tests (chapitre~\ref{ch:tests}) :

\begin{itemize}
  \item \emph{En tant que} visiteur, \emph{je veux} créer un compte \emph{afin de} publier des
        messages.
  \item \emph{En tant qu'}utilisateur, \emph{je veux} publier un message avec une image
        \emph{afin de} partager du contenu.
  \item \emph{En tant qu'}utilisateur, \emph{je veux} répondre à un message \emph{afin de}
        participer à une conversation.
  \item \emph{En tant qu'}administrateur, \emph{je veux} supprimer n'importe quel message
        \emph{afin de} modérer la plateforme.
\end{itemize}

\section{Contraintes réglementaires et qualité}

\paragraph{RGPD.} L'application traite des données à caractère personnel (adresse email, mot de
passe, contenus publiés). Le mot de passe n'est jamais stocké en clair : il est haché
(\texttt{bcrypt}). La collecte est minimale (principe de minimisation). Des mentions légales et
une information sur la finalité du traitement doivent accompagner la mise en production.

\paragraph{RGAA (accessibilité).} Les interfaces visent un socle d'accessibilité : attributs
\texttt{alt} sur les images, contrastes suffisants, libellés de formulaire, navigation au
clavier et attributs \texttt{aria-*} sur les éléments interactifs (par exemple le champ de
saisie de message porte \texttt{role="textbox"} et \texttt{aria-label}).

\paragraph{Éco-conception.} Plusieurs choix limitent l'empreinte : minification des assets en
production, pagination du fil (\texttt{LIMIT 20}), et suppression d'une requête N+1 réduisant le
nombre d'allers-retours vers la base (section~\ref{sec:acces}).

